% Source: https://artofproblemsolving.com/community/c3h1579342_helpcones
% Stable identity: aops-topic-1579342-q01
Put the apex at the top and cut the cone by horizontal planes into $n$ layers of thickness $\dfrac {h}{n}$. At depth $y$ below the apex the cross-section is a disc of radius $\dfrac{ry}{h}$ (similar triangles), so in the $k$-th layer from the top, between depths $\dfrac{(k-1)h}{n}$ and $\dfrac{kh}{n}$, the cross-sections have radius between $\dfrac{r(k-1)}{n}$ and $\dfrac{rk}{n}$.

\textbf{Outer stack.} The cylinder of radius $\dfrac{rk}{n}$ and height $\dfrac {h}{n}$ contains the $k$-th layer. These $n$ cylinders together contain the cone, so
\[
V\le\sum_{k=1}^{n}\pi\left(\dfrac{rk}{n}\right)^2\dfrac {h}{n}=\dfrac{\pi r^2h}{n^3}\sum_{k=1}^{n}k^2=\pi r^2h\cdot\dfrac{n(n+1)(2n+1)}{6n^3}.
\]
\textbf{Inner stack.} Likewise the cylinders of radius $\dfrac{r(k-1)}{n}$ lie inside the layers, so
\[
V\ge\dfrac{\pi r^2h}{n^3}\sum_{k=1}^{n}(k-1)^2=\pi r^2h\cdot\dfrac{(n-1)n(2n-1)}{6n^3}.
\]

\begin{asy}
usepackage("amsmath");
size(5cm);
defaultpen(fontsize(10pt));
real r=3,h=6; int n=5;
for(int k=1;k<=n;++k){real rk=r*k/n; real y1=h-(k-1)*h/n,y0=h-k*h/n;
 filldraw((-rk,y0)--(rk,y0)--(rk,y1)--(-rk,y1)--cycle,rgb(0.85,0.9,1),blue);}
draw((-r,0)--(r,0)--(r,h)--(-r,h)--cycle,gray+dashed);
draw((-r,0)--(0,h)--(r,0),linewidth(1));
draw((-r,0)--(r,0),linewidth(1));
dot((0,h));
draw((r+0.6,0)--(r+0.6,h),Arrows(TeXHead));label("$h$",(r+0.6,h/2),E);
draw((0,-0.45)--(r,-0.45),Arrows(TeXHead));label("$r$",(r/2,-0.45),S);
label("\mbox{}",(min(currentpicture,user=true)+max(currentpicture,user=true))/2,black); // keeps the text after the figure black when LaTeX typesets the labels inline
\end{asy}

(The sketch is a vertical cross-section: the cone inside the dashed cylinder, with the outer stack for $n=5$.)

As $n$ grows, both bounds tend to $\pi r^2h\cdot\dfrac{2}{6}=\dfrac{1}{3}\pi r^2h$, because $\dfrac{n(n+1)(2n+1)}{n^3}\to2$ and $\dfrac{(n-1)n(2n-1)}{n^3}\to2$. The volume $V$ lies between them for every $n$, so
\[
V=\dfrac{1}{3}\pi r^2h,
\]
which is one third of the cylinder's volume $\pi r^2h$. $\blacksquare$

The factor $\dfrac{1}{3}$ comes from $\sum_{k=1}^n k^2\approx\dfrac{n^3}{3}$: the cross-sectional area grows like the square of the depth. (The same argument shows that every pyramid is one third of the prism with the same base and height.)
